% GENERATED FILE. Edit canonical lessons, then run npm run generate:books.
% canonical-source-hash: fnv1a64:fd5a8096d3848300
% canonical-lessons: BN-C477-jontrona, BN-C477-pora-burn, BN-C477-tonton-kora, BN-C477-khorano, BN-C477-gongano, BN-R477-review-470-477

\chapter{Burns and Aches}
\label{ch:bn-burns-and-aches}
\hlchaptermodality{477}

\begin{chapteropening}
\textbf{By the end of this chapter:} \emph{I can talk about burning, throbbing pain and limping.}

Complete the last lesson of chapter 477: I can talk about burning, throbbing pain and limping.
\end{chapteropening}

\section[\texorpdfstring{\bn{যন্ত্রণা} (jôntrônā)}{যন্ত্রণা (jôntrônā)}]{\bn{যন্ত্রণা} (jôntrônā) — agony; torment; severe pain}

\label{lesson:BN-C477-jontrona}



\begin{quote}
Before the new one: say the Bengali for to pant, then the Bengali for to groan.
\end{quote}

\subsection*{You'll want to know: \bn{যন্ত্রণা}}
\textbf{\bn{যন্ত্রণা}} — \emph{jôntrônā} — \textquotedblleft{}agony; torment; severe pain\textquotedblright{}.

\begin{quote}
\bn{দাঁতের} \bn{যন্ত্রণায়} \bn{দাদু} \bn{সারারাত} \bn{জেগে} \bn{ছিলেন।}
\end{quote}

\begin{quote}
\emph{dā̃ter jôntrônāy dādu sārārāt jege chhilen.}
\end{quote}

\begin{quote}
Grandfather lay awake all night in agony from his tooth.
\end{quote}

\bn{ব্যথা} (taught) is any pain; \bn{যন্ত্রণা} is pain that torments, of the body or of the mind.

\begin{cousinweb}
\textbf{\bn{যন্ত্রণা}} (\emph{jôntrônā}) is a tatsama, Sanskrit \emph{yantraṇā}, which Monier-Williams gives first for restraint and constraint and then for pain and anguish; Bengali keeps the pain.
\end{cousinweb}

\begin{grammarlens}[title={What you've built}]
One more word for describing what you lived and felt.
\end{grammarlens}

\subsection*{Your turn}
\begin{itemize}
\raggedright
  \item \emph{Say it:} \emph{jôntrônā}
  \item \emph{Say it:} \emph{jôntrônā}, once more
  \item \emph{Say it:} \emph{dā̃ter jôntrônāy dādu sārārāt jege chhilen.} — Grandfather lay awake all night in agony from his tooth.
  \item \emph{Recall:} say the Bengali for to pant, then the Bengali for to groan, then say \emph{jôntrônā} again
\end{itemize}

\begin{tcolorbox}[breakable,colback=teal!4,colframe=teal!35!black,title={Before you move on}]
What does \bn{যন্ত্রণা} mean? (\textquotedblleft{}Agony; torment; severe pain\textquotedblright{}.) Say it once more.
\end{tcolorbox}
\section[\texorpdfstring{\bn{পোড়া} (poṛā)}{পোড়া (poṛā)}]{\bn{পোড়া} (poṛā) — to get burnt; to burn (by itself)}

\label{lesson:BN-C477-pora-burn}



\begin{quote}
Before the new one: say the Bengali for to groan, then the Bengali for agony.
\end{quote}

\subsection*{You'll want to know: \bn{পোড়া}}
\textbf{\bn{পোড়া}} — \emph{poṛā} — \textquotedblleft{}to get burnt; to burn (by itself)\textquotedblright{}.

\begin{quote}
\bn{গরম} \bn{তেলে} \bn{মীরার} \bn{হাত} \bn{পুড়ল।}
\end{quote}

\begin{quote}
\emph{gôrôm tele mīrār hāt puṛlô.}
\end{quote}

\begin{quote}
Mira's hand got burnt by hot oil.
\end{quote}

Past (he/she …): \bn{সেই} \bn{রুটি} \bn{পুড়ল}, \emph{sei ruṭi puṛlô}.

\begin{cousinweb}
\textbf{\bn{পোড়া}} (\emph{poṛā}) is listed by CDIAL for Bengali, as burning and being burnt, under a reconstructed form (CDIAL 8397); \bn{পোড়ানো} (taught), to burn something, is its causative.
\end{cousinweb}

\begin{grammarlens}[title={What you've built}]
One more word for describing what you lived and felt.
\end{grammarlens}

\subsection*{Your turn}
\begin{itemize}
\raggedright
  \item \emph{Say it:} \emph{poṛā}
  \item \emph{Say it:} \emph{poṛā}, once more
  \item \emph{Say it:} \emph{gôrôm tele mīrār hāt puṛlô.} — Mira's hand got burnt by hot oil.
  \item \emph{Recall:} say the Bengali for to groan, then the Bengali for agony, then say \emph{poṛā} again
\end{itemize}

\begin{tcolorbox}[breakable,colback=teal!4,colframe=teal!35!black,title={Before you move on}]
What does \bn{পোড়া} mean? (\textquotedblleft{}To get burnt; to burn (by itself)\textquotedblright{}.) Say it once more.
\end{tcolorbox}
\section[\texorpdfstring{\bn{টনটন} \bn{করা} (ṭônṭôn kôrā)}{টনটন করা (ṭônṭôn kôrā)}]{\bn{টনটন} \bn{করা} (ṭônṭôn kôrā) — to throb; to ache dully}

\label{lesson:BN-C477-tonton-kora}



\begin{quote}
Before the new one: say the Bengali for agony, then the Bengali for to get burnt.
\end{quote}

\subsection*{You'll want to know: \bn{টনটন} \bn{করা}}
\textbf{\bn{টনটন} \bn{করা}} — \emph{ṭônṭôn kôrā} — \textquotedblleft{}to throb; to ache dully\textquotedblright{}.

\begin{quote}
\bn{দাঁতটা} \bn{টনটন} \bn{করছে।}
\end{quote}

\begin{quote}
\emph{dā̃tṭā ṭônṭôn kôrchhe.}
\end{quote}

\begin{quote}
My tooth is throbbing.
\end{quote}

Past (he/she …): \bn{সেই} \bn{রোগীর} \bn{মাথা} \bn{টনটন} \bn{করল}, \emph{sei rogīr māthā ṭônṭôn kôrlô}.

\begin{cousinweb}
\textbf{\bn{টনটন} \bn{করা}} (\emph{ṭônṭôn kôrā}) is \emph{ṭônṭôn}, an imitative word for a dull, throbbing ache, with \emph{kôrā}, \textquotedblleft{}to do\textquotedblright{} (CDIAL 2814).
\end{cousinweb}

\begin{grammarlens}[title={What you've built}]
One more word for describing what you lived and felt.
\end{grammarlens}

\subsection*{Your turn}
\begin{itemize}
\raggedright
  \item \emph{Say it:} \emph{ṭônṭôn kôrā}
  \item \emph{Say it:} \emph{ṭônṭôn kôrā}, once more
  \item \emph{Say it:} \emph{dā̃tṭā ṭônṭôn kôrchhe.} — My tooth is throbbing.
  \item \emph{Recall:} say the Bengali for agony, then the Bengali for to get burnt, then say \emph{ṭônṭôn kôrā} again
\end{itemize}

\begin{tcolorbox}[breakable,colback=teal!4,colframe=teal!35!black,title={Before you move on}]
What does \bn{টনটন} \bn{করা} mean? (\textquotedblleft{}To throb; to ache dully\textquotedblright{}.) Say it once more.
\end{tcolorbox}
\section[\texorpdfstring{\bn{খোঁড়ানো} (khõṛāno)}{খোঁড়ানো (khõṛāno)}]{\bn{খোঁড়ানো} (khõṛāno) — to limp}

\label{lesson:BN-C477-khorano}



\begin{quote}
Before the new one: say the Bengali for to get burnt, then the Bengali for to throb.
\end{quote}

\subsection*{You'll want to know: \bn{খোঁড়ানো}}
\textbf{\bn{খোঁড়ানো}} — \emph{khõṛāno} — \textquotedblleft{}to limp\textquotedblright{}.

\begin{quote}
\bn{চোট} \bn{পেয়ে} \bn{কুকুরটা} \bn{খোঁড়াচ্ছিল।}
\end{quote}

\begin{quote}
\emph{choṭ peye kukurṭā khõṛāchchhilô.}
\end{quote}

\begin{quote}
The dog was limping after hurting its leg.
\end{quote}

Past (he/she …): \bn{তিনি} \bn{খোঁড়ালেন}, \emph{tini khõṛālen}.

\begin{cousinweb}
\textbf{\bn{খোঁড়ানো}} (\emph{khõṛāno}) is the verb of \emph{khõṛā}, \textquotedblleft{}lame\textquotedblright{}, which CDIAL lists for Bengali (CDIAL 3941).
\end{cousinweb}

\begin{grammarlens}[title={What you've built}]
One more word for describing what you lived and felt.
\end{grammarlens}

\subsection*{Your turn}
\begin{itemize}
\raggedright
  \item \emph{Say it:} \emph{khõṛāno}
  \item \emph{Say it:} \emph{khõṛāno}, once more
  \item \emph{Say it:} \emph{choṭ peye kukurṭā khõṛāchchhilô.} — The dog was limping after hurting its leg.
  \item \emph{Recall:} say the Bengali for to get burnt, then the Bengali for to throb, then say \emph{khõṛāno} again
\end{itemize}

\begin{tcolorbox}[breakable,colback=teal!4,colframe=teal!35!black,title={Before you move on}]
What does \bn{খোঁড়ানো} mean? (\textquotedblleft{}To limp\textquotedblright{}.) Say it once more.
\end{tcolorbox}
\section[\texorpdfstring{\bn{গোঙানো} (gongāno)}{গোঙানো (gongāno)}]{\bn{গোঙানো} (gongāno) — to moan}

\label{lesson:BN-C477-gongano}



\begin{quote}
Before the new one: say the Bengali for to throb, then the Bengali for to limp.
\end{quote}

\subsection*{You'll want to know: \bn{গোঙানো}}
\textbf{\bn{গোঙানো}} — \emph{gongāno} — \textquotedblleft{}to moan\textquotedblright{}.

\begin{quote}
\bn{ব্যথায়} \bn{বুড়ো} \bn{লোকটা} \bn{গোঙাচ্ছিল।}
\end{quote}

\begin{quote}
\emph{bæthāy buṛo lokṭā gongāchchhilô.}
\end{quote}

\begin{quote}
The old man was moaning with pain.
\end{quote}

Past (he/she …): \bn{তিনি} \bn{ব্যথায়} \bn{গোঙালেন}, \emph{tini bæthāy gongālen}.

\begin{cousinweb}
\textbf{\bn{গোঙানো}} (\emph{gongāno}) is an imitative verb for a low, wordless moan; it shows the new letter \bn{ঙ}. No dictionary we use records an older form, so its history is not settled.
\end{cousinweb}

\begin{grammarlens}[title={What you've built}]
One more word for describing what you lived and felt.
\end{grammarlens}

\subsection*{Your turn}
\begin{itemize}
\raggedright
  \item \emph{Say it:} \emph{gongāno}
  \item \emph{Say it:} \emph{gongāno}, once more
  \item \emph{Say it:} \emph{bæthāy buṛo lokṭā gongāchchhilô.} — The old man was moaning with pain.
  \item \emph{Recall:} say the Bengali for to throb, then the Bengali for to limp, then say \emph{gongāno} again
\end{itemize}

\begin{tcolorbox}[breakable,colback=teal!4,colframe=teal!35!black,title={Before you move on}]
What does \bn{গোঙানো} mean? (\textquotedblleft{}To moan\textquotedblright{}.) Say it once more.
\end{tcolorbox}
\section[(dialogue)]{Looking back: from Hair and Scalp to Burns and Aches}

\label{lesson:BN-R477-review-470-477}



\begin{quote}
Say the Bengali for baldness and dandruff.
\end{quote}

\subsection*{Your turn}
Say these aloud:

\begin{itemize}
\raggedright
  \item baldness, dandruff, a louse, to fall out, a knot of hair
  \item an illness, a symptom, to suffer from, to faint, bedridden
  \item a medicine, medical treatment, a diet for the sick, a vaccine, to recover
  \item an infection, a germ, contagious, an epidemic, to be struck down
  \item asthma, diabetes, blood pressure, cancer, jaundice
  \item a wound, a sore, to sprain, to swell, to stumble
  \item stamina, hiccups, a burp, to pant, to groan
  \item agony, to get burnt, to throb, to limp, to moan
\end{itemize}

\begin{tcolorbox}[breakable,colback=teal!4,colframe=teal!35!black,title={Before you move on}]
Baldness? (\textbf{\bn{টাক}}, \emph{ṭāk}.) Dandruff? (\textbf{\bn{খুশকি}}, \emph{khushki}.) A louse? (\textbf{\bn{উকুন}}, \emph{ukun}.) To fall out? (\textbf{\bn{ঝরা}}, \emph{jhôrā}.) A knot of hair? (\textbf{\bn{খোঁপা}}, \emph{khõpā}.) An illness? (\textbf{\bn{অসুখ}}, \emph{ôsukh}.) A symptom? (\textbf{\bn{লক্ষণ}}, \emph{lôkkhôn}.) To suffer from? (\textbf{\bn{ভোগা}}, \emph{bhogā}.) To faint? (\textbf{\bn{অজ্ঞান} \bn{হওয়া}}, \emph{ôgyān hôoyā}.) Bedridden? (\textbf{\bn{শয্যাশায়ী}}, \emph{shôjyāshāyī}.) A medicine? (\textbf{\bn{ওষুধ}}, \emph{oshudh}.) Medical treatment? (\textbf{\bn{চিকিৎসা}}, \emph{chikitsā}.) A diet for the sick? (\textbf{\bn{পথ্য}}, \emph{pôthyô}.) A vaccine? (\textbf{\bn{টিকা}}, \emph{ṭikā}.) To recover? (\textbf{\bn{সেরে} \bn{ওঠা}}, \emph{sere oṭhā}.) An infection? (\textbf{\bn{সংক্রমণ}}, \emph{sôngkrômôn}.) A germ? (\textbf{\bn{জীবাণু}}, \emph{jībānu}.) Contagious? (\textbf{\bn{ছোঁয়াচে}}, \emph{chhõyāche}.) An epidemic? (\textbf{\bn{মহামারি}}, \emph{môhāmāri}.) To be struck down? (\textbf{\bn{আক্রান্ত} \bn{হওয়া}}, \emph{ākrāntô hôoyā}.) Asthma? (\textbf{\bn{হাঁপানি}}, \emph{hā̃pāni}.) Diabetes? (\textbf{\bn{ডায়াবেটিস}}, \emph{ḍāyābeṭis}.) Blood pressure? (\textbf{\bn{রক্তচাপ}}, \emph{rôktôchāp}.) Cancer? (\textbf{\bn{ক্যানসার}}, \emph{kænsār}.) Jaundice? (\textbf{\bn{জন্ডিস}}, \emph{jônḍis}.) A wound? (\textbf{\bn{ক্ষত}}, \emph{khôtô}.) A sore? (\textbf{\bn{ঘা}}, \emph{ghā}.) To sprain? (\textbf{\bn{মচকানো}}, \emph{môchkāno}.) To swell? (\textbf{\bn{ফোলা}}, \emph{pholā}.) To stumble? (\textbf{\bn{হোঁচট} \bn{খাওয়া}}, \emph{hõchôṭ khāoyā}.) Stamina? (\textbf{\bn{দম}}, \emph{dôm}.) Hiccups? (\textbf{\bn{হেঁচকি}}, \emph{hẽchki}.) A burp? (\textbf{\bn{ঢেকুর}}, \emph{ḍhekur}.) To pant? (\textbf{\bn{হাঁপানো}}, \emph{hā̃pāno}.) To groan? (\textbf{\bn{কাতরানো}}, \emph{kātrāno}.) Agony? (\textbf{\bn{যন্ত্রণা}}, \emph{jôntrônā}.) To get burnt? (\textbf{\bn{পোড়া}}, \emph{poṛā}.) To throb? (\textbf{\bn{টনটন} \bn{করা}}, \emph{ṭônṭôn kôrā}.) To limp? (\textbf{\bn{খোঁড়ানো}}, \emph{khõṛāno}.) To moan? (\textbf{\bn{গোঙানো}}, \emph{gongāno}.)
\end{tcolorbox}
